\documentclass[10pt]{article}
\usepackage[a4paper,margin=25mm]{geometry}
\usepackage[T1]{fontenc}
\usepackage{lmodern}
\usepackage{amsmath,amssymb,booktabs,array}
\usepackage[colorlinks=true,urlcolor=blue,citecolor=blue]{hyperref}
\usepackage{listings}
\lstset{basicstyle=\ttfamily\footnotesize,breaklines=true,columns=fullflexible,keepspaces=true,frame=single}
\title{ChoreoFix: Enforcing Multi-Agent Coordination from Global Protocols}
\author{Hayato Ohhashi}
\date{15 October 2026\\Research draft, comparative evaluation incomplete}
\begin{document}
\maketitle
\begin{abstract}
Agents can choose how to work, but shared effects still need clear rules.
We study whether a global protocol can serve as the source of those rules
without a separate application state machine. ChoreoFix uses Hibana to project
Rust global protocols into local communication programs. Handwritten async
roles drive their endpoints before carrying out domain effects. A bounded WASI
shell lets agents use work files and scripts. Every file operation goes through
ChoreoFS. The shell cannot grant a task step by choosing a pathname or writing
an approval file.
The current direct boundary has been tested on two TraceFix tasks.
One three-participant CI/CD trial and three three-participant review trials
used GPT-6 Luna with xhigh reasoning. The review trials covered zero, one,
and two requested revisions. All four completed their protocols and original
simulator work without recorded resource violations. Independent audits check
saved requests, effects, replies, model identity, and inclusive token counts.
These are unscored integration observations, not a matched baseline comparison.
The planned comparison retains all 48 tasks and four original runtime
conditions. No task has passed the complete benchmark acceptance gate and no
scored episode has started.
\end{abstract}

\section{Problem and scope}
A group of agents is a distributed program. An agent may prepare a change,
review a result, or choose another tool. It must still wait for the evidence or
resource needed by a shared operation. Putting an order in a prompt does not
make the effect depend on that order.

A runtime can check the order. The question here is how to define and maintain
that check. We write coordination as global protocols, project the local
communication programs, and let affine endpoints track progress. We keep
application computation in ordinary Rust. We do not maintain an application
program counter, completion flag, phase enum, or resource-owner map alongside
the endpoints.

The name ChoreoFix uses ``fix'' to mean fixing the permitted coordination
relation. It does not claim counterexample-guided automatic protocol repair.
Our research question is:
\begin{quote}
Can agents choose local work methods while projected global protocols guard
shared effects, retries, and independent work without a second application
progress authority?
\end{quote}

Three kinds of evidence are needed. A task comparison must measure safety and
completion on the original 48 tasks. A boundary experiment must show that
several legal command sequences work and that forbidden effects stop before
execution. A maintenance study must measure changes to declarations, localside
code, and authorization state. The present draft contains implementation and
integration evidence. The full task comparison and maintenance study are open.

\section{Background}
TraceFix synthesizes coordination protocols, checks PlusCal/TLA+ models, and
uses counterexamples to repair them. Its runtime study compares four ways to
execute the work \cite{tracefix}. We retain those original conditions. In
particular, a mediator already drives coordination. Runtime enforcement alone
is therefore not our claimed contribution. We test a different implementation
boundary and account for its limits.

Hibana takes direct inspiration from affine multiparty session types and
MultiCrusty \cite{ampst}. A global protocol describes communication among
roles. Projection derives each role's communication obligations. Hibana keeps
the resulting local program as data. An endpoint consumes permitted operations
without changing its Rust type after each step. Projection does not generate
the participant function. The function is handwritten, and its communication
is checked during execution.

\section{Execution boundary}
\subsection{Global and localside}
The implementation has three visible parts. Each task's \texttt{global.rs}
defines its workflow and handoffs. Its \texttt{localside/} functions show the
corresponding endpoint operations. The \texttt{imp/} code carries bytes and
model or simulator RPCs. It does not decide the next task step.

Sequential composition represents dependencies. Parallel composition retains
independent work. Routes represent alternatives, declared resolvers use actual
domain results, and repetition represents another attempt. These constructors
are used where the task needs them. A terminal test is not turned into a retry
loop merely to use a route.

A successful local step authorizes only the next operation represented by that
endpoint. It does not prove the quality of the agent's computation. The domain
operation is called by trusted localside code after the endpoint admits the
corresponding step. A failed operation cannot silently select another branch.
Task recovery follows an explicit route. A transport or session fault ends the
run and is not treated as a successful domain retry.

The review task, 6E, has two developers, a reviewer, and a shared repository.
The implementation assigns WASI roles, feature owners, reviewer services, and
one repository owner within the same global. The reviewer is one model and
one guest. That guest directly owns a caller endpoint for each service and a
private-I/O endpoint. Its two service owners do not add model participants.
Each feature can proceed while the other waits for its reviewer. A revision
uses the actual review result at an explicit resolver point:
\begin{lstlisting}
let feature_a = g::seq(
    implement_a,
    g::route(accepted_a, revise_a).resolve::<DECIDE_A>(),
).roll();
\end{lstlisting}
The complete global defines implementation requests, code-ready messages,
review requests and replies, repository operations, and completion. It composes
both features with \texttt{g::par}. The accepted arm admits merge and the
revision arm returns feedback before another implementation. Each corresponding
localside writes its \texttt{send}, \texttt{recv}, and \texttt{offer} operations
directly. For example, the developer consumes actual feedback:
\begin{lstlisting}
let feedback = endpoint.offer().await?;
match feedback.label() {
    30 => {
        let value = feedback.recv::<Approved<'_>>().await?;
        // Return this review result, then serve the merge.
    }
    31 => {
        let value = feedback.recv::<Revise<'_>>().await?;
        // Return this review result, then serve another implementation.
    }
    _ => return Err(...),
}
\end{lstlisting}
This is an excerpt, not a complete function. The source contains the actual
serialization and endpoint replies. Labels identify projected alternatives.
No approved flag or stored next-operation enum selects the continuation.
The resolver input holds the original review result, not task progress.

\subsection{Resources and notifications}
The sole repository owner receives projected requests before applying the
original implementation or merge operation. It owns the service object bytes
and awaits each domain operation before serving another. No separate lock
owner table grants access. This relies on the trusted single materializing
owner. Endpoint conformance does not by itself prove arbitrary Rust side
effects mutually exclusive.

Implementation and code-ready receipt precede review. The actual result reaches
the developer through the review-feedback channel. Approval must precede merge.
The global also carries the result to the guest. The trace auditor checks these
channels, effect intervals, and both completed merges separately from the
runtime's completion message.

The 16E CI/CD task connects its three WASI roles directly to the developer,
tester, deployer, and two resource roles in one global. Resource acquire, grant,
return, and acknowledgement are messages. Build completion precedes testing,
test completion precedes deployment, and deployment completion notifies both
other participants. Its original simulator has no injected domain failure.
A retry loop is therefore not added to this task.

Earlier codebase and writing prototypes connect task and resource sessions in
Rust brokers. Even when both sides use endpoints, that connection is outside
a single task global. Those prototypes and their tests are retained as component
evidence. They do not meet the target architecture and are not accepted task
ports. Porting must preserve disjoint work, original resource counts, and
terminal outcomes while replacing these connections. A serialized substitute
is not sufficient for tasks that allow disjoint operations.

\subsection{WASI and ChoreoFS}
The model sees one shell tool. Commands enter a real Wasm guest under
\texttt{hibana-wasip1-runtime}. Its pending WASI import is sent through a
projected ABI request/reply protocol. A matching typed completion resumes the
guest. ChoreoFS supplies path, object, fd, and I/O facts to the answering role.
Work files, scripts, terminal objects, and tool request/reply objects all use
this path. There is no guest-side substitute filesystem or host filesystem
fallback.

The current guest is the independent Rust \texttt{std} program
\texttt{hibana-wash}. It has no external Rust dependencies. It supports quoted
arguments, variables, conditions, loops, pipelines, redirection, calculations,
and saved scripts. Eight private files have 8,192-byte capacities. A script has
a 4,096-byte UTF-8 limit. Each guest starts with 64 KiB committed memory and can
request growth through the projected protocol up to 1 MiB of supplied backing.
These are explicit experimental limits, not an unrestricted POSIX environment.

The model receives object paths and the original domain schemas. It chooses
its commands. Writing JSON arguments to a domain object and closing it requests
that operation. The actual path is:
\begin{quote}
WASI endpoint $\rightarrow$ projected request $\rightarrow$ owner endpoint
$\rightarrow$ ChoreoFS operation $\rightarrow$ projected reply.
\end{quote}
The task global includes those edges. There is no unprojected request queue
connecting an ABI graph to a separate task graph. Private work uses a parallel
projected I/O fragment and cannot supply a missing task permission.

Object binding and argument validation require the target to match the chosen
service. A rejected argument follows an explicit resolver branch and can be
corrected without restarting the endpoint. ChoreoFS handles object operations,
not task authority. An open fd or an approval file cannot grant a merge.
The older limited-command guest and generic \texttt{/tools/call} bridge remain
only as historical component evidence.

\subsection{Trusted code and claim boundary}
The boundary assumes the supplied tool interface is the only path to the
protected effects. Model roles cannot call native tools or choose another
participant's identity. It also trusts the localside effect adapter, the
runtime, codecs, transport, executor, and simulator integration. The simulator
independently records actual domain effects and resource violations.

Endpoint conformance is not a proof of arbitrary program termination, review
quality, or artifact correctness. Existing Hibana and WASI proofs cover stated
models and assumptions. We do not claim a refinement proof for every line of
this application or a whole-system distributed deadlock theorem.

\section{Evaluation design}
\subsection{Original task comparison}
The selected matrix has 48 tasks, three failure depths, three repetitions,
and five conditions, for 2,160 scored episodes. The conditions are the paper's
Topology-monitored, Mediator-enforced, Prompt-only, Chat-only, and Hibana.
All participants use \texttt{gpt-6-luna} with xhigh reasoning. The episode limit
is 600 seconds and each participant has a 200-tool-call budget.

The four TraceFix runtimes remain at public revision
\href{https://github.com/Sensing-And-Reasoning/TraceFix/tree/87b01a2c0ae53c5045c9ce3d43cd0235b1831275}{\texttt{87b01a2}}. Their coordination logic is
unchanged. Benchmark descriptions and tool metadata at recorded revision
\href{https://github.com/Sensing-And-Reasoning/TraceFix/tree/36b4bd108fcd74f7efb17613aebbe3fdf55499a9}{\texttt{36b4bd1}} are byte-identical to that
release. A separate condition uses the unchanged latest monitoring runtime at
\texttt{36b4bd1}, adding 432 paired episodes across the same tasks, failure
depths, and repetitions. The total plan has 2,592 episodes. Its existing state,
monitor, distributed-coordination, and corrective-guidance tests passed in 94
cases. This is source qualification, not a measured agent comparison.
Our model and budgets differ from the original paper, so its percentages are not
matched controls for our runs.

Every port must retain the original roles, resources, dependencies, failure
paths, and required independence. All 825 checklist items remain in the
manifest. A passing compilation or generated TLC model is not by itself a
fidelity verdict. Task prototypes exist, but none has passed the
full independent benchmark acceptance gate. No scored episode has started.

\subsection{Outcomes and accounting}
Primary outcomes are full and partial completion and unauthorized committed
effects. We distinguish pre-effect rejection, domain failure, incomplete work,
timeout, protocol fault, and harness fault. We also record resource overlap,
latency, tool calls, and inclusive tokens. Cached inputs remain part of token
use. Reasoning tokens are not added twice.

The comparison must use paired task, failure-depth, and repetition identities.
Confidence intervals and paired contrasts will be computed only after scored
runs exist. Setup failures and unscored integration runs remain visible and
outside the scored denominators. A rejected effect is not a completed task.

\subsection{Separate shell experiment}
The shell changes the original tool interface. Its integration results are
therefore separate from the 2,160-run comparison. TraceFix keeps the original
interface. A causal comparison of shell coordination would require the same
guest, objects, prompts, and budgets in both arms.

The model receives the original task description, its assigned input data, and
the domain tool schemas. It receives no fixed command recipe. Boundary tests
also exercise scripts and several legal file-operation sequences. A private
``approved'' file is deliberately written before an early shared operation.
The task endpoint still rejects the operation before the simulator receives
an execute request.

\section{Evidence to date}
\subsection{Source and core qualification}
The current consumer uses Hibana package version \texttt{0.9.9} with an exact Git
patch to \texttt{c2001ec06944d7f3b299e16b5b3efb5b0ba0366a}. This is an unreleased
repair, not a new registry version. Both direct and WASI-runtime dependency
paths select that core. The WASI runtime is pinned to
\texttt{4ebdad1443cf49488eae924afe528f67f5a5953b} and wash to
\texttt{12e137aba2cea74f353c198a8d7c92afc977627c}.

Developing the 6E boundary exposed projection ambiguity at an elastic repetition
boundary and runtime problems involving parallel resolver visits and retained
send-selected choices. Minimal native regressions reproduce the affected cases.
The repair retains valid independent parallel repetitions. Local Rust tests
passed 883 cases with 12 explicit exclusions. The saved verification includes
Lean, Z3, focused Miri and Kani checks, embedded compilation, and Clippy.
These checks cover their recorded models and cases. They do not prove every
application trace safe. The core's separate full CI run 38734656111 also passed
both final-form and Kani at that exact commit.

A QUIC consumer selecting the same exact core passed 593 default and 598 HTTP/3
local tests. External run 38735739939 passed all 44 registered case/direction
cells at QUIC revision 44d291a in one run. Saved phase reports, verdicts,
expected cells, source-pair trees, and core/TLS pins were checked independently.
Reference self-tests were omitted. The retained evidence discloses runner and
capture-shutdown patches. This is integration evidence for the repair, not a
whole-system proof or a performance comparison.

\subsection{Direct WASI model trials}
Table~\ref{tab:preflight} contains four unscored trials using the direct projected
boundary. Earlier limited-command and cross-graph bridge trials remain in their
separate ledgers. They are not evidence that the present architecture was
already complete. All participants in this table used the requested and returned
Luna model and xhigh effort. The original simulator and task inputs were kept.
\begin{table}[h]
\centering
\caption{Unscored direct WASI trials. No baseline contrast is available.}
\label{tab:preflight}
\begin{tabular}{lrrrrrr}
\toprule
Task & Depth & Roles & Shell calls & Domain calls & Seconds & Tokens \\
\midrule
16E & 0 & 3 & 11 & 3 & 27.127 & 220,215 \\
6E & 0 & 3 & 20 & 6 & 59.481 & 361,915 \\
6E & 1 & 3 & 23 & 8 & 45.387 & 412,438 \\
6E & 2 & 3 & 22 & 10 & 46.913 & 397,336 \\
\bottomrule
\end{tabular}
\end{table}

The 16E observation uses the pre-repair published core and is retained at its
own source checkpoint. All four required cross-role notifications and three
domain effects were audited. The three 6E observations use consumer 4a9937f and
core c2001ec. Each completed both merges with no recorded simulator violation.
The source snapshot was frozen before model calls and is identical across the
three review trials. Binaries, compiler details, prompts, raw model events, raw
carrier events, and inclusive usage were retained. Source hashes alone are not
used as evidence that a particular model turn finished.

For 6E the original simulator fails the first depth review calls across all
features. At depth 2 the model chose review order A, B, A, B. Each feature was
revised once. Deterministic boundary cases instead exercise A first and B first
with zero, one, or two revisions concentrated on the first feature. These cover
different observed schedules without requiring the model to follow a recipe.
The three model trials used 1,171,689 inclusive tokens in total. This is a cost
observation, not evidence of a token saving.

The read-only auditor replays the original simulator and checks every actual
close before its domain effect, nonoverlapping shared-resource intervals,
code-ready delivery, review-feedback contents, revised implementation after
rejection, merge after approval, and complete frame consumption. It also checks
raw model identity and turn completion, shell command correspondence, source
and binary hashes, and inclusive token arithmetic. Altered copies remove or
forge required requests and notifications, alter feedback, or overlap resource
intervals. They must fail even after the event-file hash is recomputed.

Later boundary tests found a caller-side limitation in that model checkpoint.
After a duplicate request to a completed review service, its next branch could
be observed by an offer intended for the other service. An offer is not scoped
by the previous request. All six added cases failed in the old application.
The revised global gives each service a separate caller role, directly driven
by the same original guest. All six repaired cases and simulator replays passed,
including chunked writes and duplicate requests followed by work on the other
feature. No application progress state or core change was needed. The earlier
model observations retain their original source and do not qualify this revised
source by themselves.

\subsection{What these trials do not establish}
These are integration observations from development, not random samples or
paired baseline trials. They cannot estimate a comparative completion rate,
latency advantage, or maintenance saving. Earlier failed and interrupted
attempts remain in the repository records. A successful replacement is not
used to erase a failed attempt.

The shell changes the domain-call interface. In 6E a developer's implementation
reply waits for review and includes its actual result. Canonical object targets
are checked before execution. The unchanged baseline retains its own argument
matching and original tool replies. These differences must be controlled in a
causal shell comparison or reported as a separate interface experiment.

The original tools simulate work and do not consume an agent's source patch.
They can measure coordination, not whether one reasoning strategy produces a
better software repair. Private shell computation is flexible within the
stated profile. It does not make this a general operating-system sandbox or a
proof of unrestricted agent freedom. No complete 48-task fidelity acceptance,
scored comparison, or maintenance measurement is reported here.

\section{Maintenance study and remaining work}
A maintenance study should change required evidence, retry rules, or a
participant relation while keeping a fixed task variant in both implementations.
For each change we will record edited declarations, localside functions,
handwritten authorization fields, tests, and failure of stale implementations.
Generated files and generic engine code must be counted separately from manual
edits. TraceFix may also derive execution data from one source. A claim of lower
maintenance cost must be measured rather than inferred from filenames.

The main blockers to a completed empirical paper are completion of all task
ports and independent fidelity acceptance of all 48, a runnable and accepted
baseline corpus, the 2,160 main and 432 supplemental scored episodes, analysis
with uncertainty, and the
maintenance measurements. The shell profile also limits the kinds of work an
agent can carry out. The original simulator checks coordination outcomes and
does not establish the quality of arbitrary generated source code.

\section{Conclusion}
ChoreoFix provides an executable way to keep workflow progress in projected
endpoints while agents choose commands in a ChoreoFS-backed WASI shell.
The present evidence covers resource handoffs, explicit test retries,
independent work, and rejection before a forbidden task effect. It supports
further evaluation. It does not yet support a comparative performance claim or
a submission-ready empirical conclusion.

\begin{thebibliography}{9}
\bibitem{tracefix} Shuren Xia, Qiwei Li, Taqiya Ehsan, and Jorge Ortiz.
\emph{TraceFix: Repairing Agent Coordination Protocols with TLA+ Counterexamples}.
2026. \url{https://arxiv.org/abs/2605.07935}.
\bibitem{ampst} Nicolas Lagaillardie, Rumyana Neykova, and Nobuko Yoshida.
\emph{Stay Safe under Panic: Affine Rust Programming with Multiparty Session Types}.
ECOOP 2022. \url{https://arxiv.org/abs/2204.13464}.
\end{thebibliography}
\end{document}
